\begin{abstract}
% This report presents a unified technical account of the StepAudio 2.5 model family, covering the shared audio-language foundation and its ASR, TTS, and Realtime specializations. We argue that StepAudio 2.5 is best understood through a shared sequence modeling interface over text and audio, established by staged multimodal pretraining and subsequently specialized through task-specific data, objectives, and decoding strategies. Within this framework, the ASR branch leverages the acoustic determinism of transcription to support verifiable multi-token decoding, yielding a favorable quality-latency trade-off. The TTS branch casts speech generation as semantic-to-audio alignment under natural-language control, and improves controllability and expressivity through context-rich supervision and reinforcement learning. The Realtime branch targets low-latency spoken interaction with persona consistency, paralinguistic perception, and response generation conditioned on conversational state. Taken together, these systems illustrate a broader design principle: once text and audio are embedded in a sufficiently strong shared representation space, recognition, synthesis, and realtime interaction can be built as different specializations of a common foundation model.

% 叙述的更满
% Unified audio-language modeling has emerged as a prominent trend in modern speech systems, promising shared representations and the broader reasoning capabilities of large language models for speech tasks. However, existing unified models do not consistently match specialized systems across automatic speech recognition, text-to-speech synthesis, and realtime spoken interaction, because these capabilities require substantially different deployment objectives that are difficult to satisfy simultaneously within a single foundation. Achieving specialized-level performance across all three capabilities within a single audio-language system remains a central challenge. This report presents StepAudio 2.5, a unified audio-language model family designed for production-grade ASR, TTS, and realtime spoken interaction. The system is built on a shared multimodal foundation established through large-scale audio-language pretraining, while task-tailored reinforcement learning serves as the key post-training mechanism for aligning the shared foundation with the distinct deployment requirements of each capability. Each branch instantiates this principle differently: the ASR branch combines autoregressive recognition with verifiable multi-token decoding for accurate and efficient transcription; the TTS branch uses context-rich supervision with reinforcement learning over instruction-following and human-preference signals for controllable, expressive synthesis under natural-language control; and the Realtime branch is aligned for low-latency dialogue with persona consistency, paralinguistic sensitivity, and conversational coherence through generative reward modeling. These specializations demonstrate that, with task-tailored alignment, a single audio-language foundation can achieve state-of-the-art performance across ASR, TTS, and realtime spoken interaction, combining the breadth of a general-purpose system with the depth of specialized ones.

% 更开放的描述
% Unified audio-language modeling has emerged as a prominent trend in modern speech systems, promising shared representations and the broader reasoning capabilities of large language models for speech tasks. However, existing unified models do not consistently match specialized systems across automatic speech recognition (ASR), text-to-speech synthesis (TTS), and realtime spoken interaction (Realtime), whose deployment objectives are difficult to satisfy simultaneously within a single foundation. Bridging the breadth of a unified foundation with the depth of specialized systems remains an open direction for audio-language modeling. This report presents StepAudio 2.5, a unified audio-language model family covering ASR, TTS, and Realtime. The system is built on a shared multimodal foundation established through large-scale audio-language pretraining, and task-tailored reinforcement learning plays a key role in post-training, aligning it with the distinct deployment requirements of each capability. Across the three branches, this takes different forms: the ASR branch uses verifiable multi-token decoding for accurate and efficient transcription; the TTS branch combines context-rich supervision with preference-based reinforcement learning for controllable, expressive synthesis under natural-language control; and the Realtime branch is aligned through generative reward modeling for low-latency dialogue with persona consistency, paralinguistic sensitivity, and conversational coherence. On standard benchmarks for the three capabilities, StepAudio 2.5 achieves state-of-the-art results across ASR, TTS, and Realtime, outperforming both leading unified and specialized systems. These results represent a substantive step toward unifying speech understanding, generation, and interaction within a single audio-language foundation.

% Unified audio-language modeling has emerged as a prominent trend in modern speech systems, promising to bring the reasoning capabilities of large language models to auditory tasks. However, existing unified foundations often struggle to match the depth of specialized systems across automatic speech recognition (ASR), text-to-speech synthesis (TTS), and realtime spoken interaction. Bridging this gap remains an open challenge. This report presents StepAudio 2.5, a unified audio-language foundation model that matches or exceeds specialized systems across all three capabilities. \textbf{Rather than treating these tasks as architecturally distinct, we operate on the premise that once text and audio share a multimodal representational space, task specialization becomes a matter of operational regimes: data construction, optimization targets, and decoding constraints.} Guided by this insight, we leverage task-tailored reinforcement learning and capability-specific post-training to shape a shared backbone into three distinct operational modes. Concretely, the ASR branch advances transcription efficiency via verifiable multi-token decoding; the TTS branch achieves controllable, expressive synthesis through preference-based reinforcement learning and context-rich supervision; and the Realtime branch realizes low-latency, persona-consistent dialogue via generative reward modeling. On standard benchmarks, StepAudio 2.5 achieves state-of-the-art results across ASR, TTS, and Realtime, demonstrating that a singular audio-language foundation can successfully internalize the distinct deployment objectives of speech understanding, generation, and live interaction.

Unified audio-language modeling has emerged as a prominent trend in modern speech systems, promising to bring the reasoning capabilities of large language models to auditory tasks. However, existing unified foundations often struggle to match the depth of specialized systems across automatic speech recognition (ASR), text-to-speech synthesis (TTS), and realtime spoken interaction. Bridging this gap remains an open challenge. This report presents StepAudio 2.5, a unified audio-language foundation model that matches or exceeds specialized systems across all three capabilities. \textbf{Rather than treating these tasks as architecturally distinct, we operate on the premise that once text and audio share a multimodal representational space, task specialization becomes a matter of operational regimes: data construction, optimization targets, and decoding constraints.} Guided by this insight, we advance the post-training paradigm from standard supervised learning to task-tailored Reinforcement Learning from Human Feedback (RLHF), using it as the primary mechanism to define complex optimization targets. We leverage this RLHF-centric alignment, alongside specialized decoding, to shape a shared backbone into three distinct operational modes. Concretely, the ASR branch advances transcription efficiency via verifiable multi-token decoding; the TTS branch achieves controllable, expressive synthesis through preference-based RLHF and context-rich supervision; and the Realtime branch realizes low-latency, persona-consistent dialogue via generative reward modeling within an RLHF framework. On standard benchmarks, StepAudio 2.5 achieves state-of-the-art results across ASR, TTS, and Realtime, demonstrating that a singular audio-language foundation can successfully internalize the distinct deployment objectives of speech understanding, generation, and live interaction.


\end{abstract}